% Generated draft: requires mathematical and source review.
\documentclass[11pt]{article}
\usepackage[margin=25mm]{geometry}
\usepackage{fontspec}
\setmainfont{Latin Modern Roman}
\usepackage{amsmath,amssymb,mathtools,unicode-math}
\setmathfont{Latin Modern Math}
\usepackage{enumitem,needspace,xurl,hyperref}
\setlength{\parindent}{0pt}
\setlength{\parskip}{5pt}
\setlength{\emergencystretch}{3em}
\newcommand{\N}{\mathbb{N}}
\newcommand{\Z}{\mathbb{Z}}
\newcommand{\Q}{\mathbb{Q}}
\newcommand{\R}{\mathbb{R}}
\newcommand{\C}{\mathbb{C}}
\newcommand{\F}{\mathbb{F}}
\newcommand{\A}{\mathbb{A}}
\newcommand{\PP}{\mathbb{P}}
\newcommand{\E}{\mathbb{E}}
\newcommand{\eps}{\varepsilon}
\newcommand{\dd}{\,\mathrm{d}}
\begin{document}
\section*{UnsolvedMath ID: 30005414}
\textbf{Source problem code:} catalogue entry
\subsection{Definitions and mathematical statement}
Let \(R\) be a commutative Noetherian ring with identity and prime characteristic \(p>0\). Write \(R^{\circ}=R\setminus\bigcup_{\mathfrak p\in\operatorname{Min}(R)}\mathfrak p\). Excellence means that \(R\) is universally catenary, every finitely generated \(R\)-algebra has open regular locus, and, for every prime \(\mathfrak p\), the completion map \(R_{\mathfrak p}\to\widehat{R_{\mathfrak p}}\) has geometrically regular fibers.

For \(e\geq 0\), put \(q=p^e\). Let \(R^{(e)}\) denote the bimodule whose underlying left \(R\)-module is \(R\), with right action \(a\cdot r=ar^q\), and set \(F_R^e(M)=R^{(e)}\otimes_R M\). For an inclusion \(N\subseteq M\) of \(R\)-modules, let \(N_M^{[q]}\) be the image of \(F_R^e(N)\to F_R^e(M)\). Its tight closure is
\[
N_M^*=\left\{x\in M:\ \exists c\in R^{\circ}\ \exists e_0\geq0\ \forall e\geq e_0,\quad c(1\otimes x)\in N_M^{[p^e]}\right\}.
\]
In particular, for an ideal \(I\subseteq R\), writing \(I^{[q]}=(a^q:a\in I)\), this gives
\[
I^*=\left\{x\in R:\ \exists c\in R^{\circ}\ \exists e_0\geq0\ \forall e\geq e_0,\quad cx^{p^e}\in I^{[p^e]}\right\}.
\]
The finitistic tight closure is
\[
N_M^{*\mathrm{fg}}=
\bigcup_{\substack{L\subseteq M\\L\text{ finitely generated}}}
\operatorname{im}\bigl((N\cap L)_L^*\longrightarrow M\bigr).
\]
The ring \(R\) is \emph{weakly \(F\)-regular} if every ideal is tightly closed; equivalently, \(N_M^{*\mathrm{fg}}=N\) for every inclusion of modules. Following the module-theoretic convention indicated by the supplied excerpt, \(R\) is \emph{strongly \(F\)-regular} if \(N_M^*=N\) for every inclusion \(N\subseteq M\) of arbitrary \(R\)-modules. For comparison with the supplied background, \(R\) is \emph{\(F\)-regular} if \(R_{\mathfrak p}\) is weakly \(F\)-regular for every \(\mathfrak p\in\operatorname{Spec}(R)\).

\textbf{Conjecture 1 (the weak implies strong conjecture).} For every excellent ring \(R\) of prime characteristic \(p>0\), does
\[
\bigl(\forall I\subseteq R\text{ an ideal},\ I^*=I\bigr)
\quad\Longrightarrow\quad
\bigl(\forall R\text{-modules }N\subseteq M,\ N_M^*=N\bigr)
\]
hold?

The quantification over modules on the right includes modules that are not finitely generated. No locality, domain, finite-type, or \(F\)-finiteness hypothesis is imposed on the conjecture. The source extraction damages the notation distinguishing finitistic and unrestricted tight closure; the definitions above reconstruct that distinction and require mathematical review against the report. In particular, a definition restricted to finite Frobenius maps has not been substituted for the supplied general-ring formulation.
\subsection{Short English statement}
Must every excellent weakly \(F\)-regular ring of prime characteristic \(p>0\) be strongly \(F\)-regular? In the module-theoretic formulation used here, the question asks whether tight closure being trivial for all ideals forces it to be trivial for every inclusion of arbitrary modules. The supplied record labels the problem partially solved; that status and the reconstructed definitions require independent review.
\subsection{Sources}
\begin{itemize}
\item Supplied local UnsolvedMath registry record, permanent ID \(30005414\), source problem code \texttt{catalogue entry}, title \emph{Weak and Strong F-Regularity in Excellent Rings}. The record assigns the year 2023 and supplies no individual proposer.
\item \emph{Resolutions in Local Algebra and Singularity Theory} (2023), Oberwolfach Reports, pp.\ 327--395, DOI \texttt{10.4171/owr/2023/6}: \url{https://doi.org/10.4171/owr/2023/6}. The supplied excerpt calls the target Conjecture 1, the weak implies strong conjecture. This is also the supplied literature link labeled ``original OWR conjecture.''
\item Supplied literature link labeled ``Aberbach--Polstra, dimension four'': \url{https://doi.org/10.1016/j.jalgebra.2022.04.020}.
\item Supplied literature link labeled ``arbitrary-dimensional partial progress'': \url{https://par.nsf.gov/servlets/purl/10501870}.
\item Literature assessment dated exactly \texttt{2026-08-22}: the supplied curation labels the problem partially solved and states that weak \(F\)-regularity implies strong \(F\)-regularity in several major classes, including dimensions at most four under hypotheses, while the excellent-ring conjecture remains open in full generality. This is supplied dated curation, not an independently verified literature assessment.
\item The registry labels the statement \texttt{corrected\_verified} and reports that the official Oberwolfach report was checked directly to restore definitions and remove extraction artifacts. That verification is a claim supplied by the registry; this draft does not claim an independent source visit or verification.
\end{itemize}
\end{document}
